\chapter{Methodology} \label{chap:methodology}

The methodology is based on data-driven empirical research to investigate logging practices and challenges
and guidelines from different perspectives.
It collects and combines data from software developers, academic literature and open-source software repositories.
First, logging practices and challenges encountered during software development
and maintenance are investigated through a developer study.
Second, academic literature is analysed to identify actionable logging
recommendations, literature-derived classifications and supporting empirical,
conceptual or tool-oriented evidence.
Third, logging guidance identified in open-source repositories is extracted,
classified and analysed according to recurring themes and characteristics, including
what to log, where to log, how to log and log-level verbosity.
Finally, the logging catalogue is constructed from the open-source repository
guidance and organised by themes into a consolidated set of logging guidelines.

\section{Research Design}

The research design of this thesis investigates software logging from three
perspectives: developer practices and challenges, logging guidelines
documented in academic literature and logging guidance documented in open-source
software repositories.
\\
The study is structured according to four research questions. First,
\hyperref[table:t1rq]{\textbf{RQ1}} investigates logging practices and challenges
encountered by software developers during software development and maintenance.
This research question is addressed through a developer survey.
\\
Second, \hyperref[table:t1rq]{\textbf{RQ2}} investigates logging guidelines and practices documented
in academic literature. Relevant publications are identified, selected and examined
for actionable recommendations concerning software logging.
\\
Third, \hyperref[table:t1rq]{\textbf{RQ3}} investigates logging guidance and practices documented
in open-source software repositories. Relevant repository material is collected,
prepared and analysed to identify logging recommendations and construct an
open-source guideline catalogue.
Publications that provided relevant empirical, conceptual or tool-oriented evidence 
but did not contain a recommendation were retained as supporting 
literature but were not treated as independent recommendation targets.
\\
The logging guidelines identified for \hyperref[table:t1rq]{\textbf{RQ2}} and
\hyperref[table:t1rq]{\textbf{RQ3}} remain traceable to their respective source sets.
The open-source recommendations are normalised, classified, consolidated and validated
to construct the logging guideline catalogue.
The academic recommendations establish the literature-based foundation of the
deductive codebook and are subsequently compared with the consolidated catalogue
to assess the extent of academic coverage.
\\
Finally, \hyperref[table:t1rq]{\textbf{RQ4}} investigates the themes represented
within the resulting logging guideline catalogue.
A deductive classification is first applied using the predefined logging
themes from the literature. The classified guideline content is then examined to identify
recurring concepts and additional themes across the collected guidelines.
\\
The developer study used for \hyperref[table:t1rq]{\textbf{RQ1}} remains analytically separate from the
construction of the guideline catalogue. Relationships between developer
practices and challenges and the identified guideline themes are considered later
during the discussion of the findings rather than being used to construct the themes
for \hyperref[table:t1rq]{\textbf{RQ4}}.
\\
An overview of the research design is shown in
Figure~\ref{fig:research_design}. The individual methodological stages are described
in the following sections.

\begin{figure}[H]
    \centering
    \resizebox{0.90\textwidth}{!}{%
        \begin{tikzpicture}[
            box/.style={
                rectangle,
                rounded corners=2pt,
                draw=black,
                align=center,
                text width=3.4cm,
                minimum height=1.0cm,
                inner sep=4pt,
                font=\small
            },
            resultbox/.style={
                rectangle,
                rounded corners=2pt,
                draw=black,
                align=center,
                text width=3.4cm,
                minimum height=1.0cm,
                inner sep=4pt,
                font=\small
            },
            mainarrow/.style={
                -{Latex[length=2.2mm]},
                semithick
            },
            auxarrow/.style={
                -{Latex[length=2.0mm]},
                dashed,
                thin
            },
            arrowlabel/.style={
                font=\scriptsize,
                fill=white,
                inner sep=1.5pt,
                align=center
            }
        ]

        % ------------------------------------------------------------
        % Research questions / data sources
        % ------------------------------------------------------------

        \node[box] (developer) at (-4.3,0)
        {Developer Study\\
        \hyperref[table:t1rq]{\textbf{RQ1}}};

        \node[box] (literature) at (0,0)
        {Academic Literature\\
        Collection\\
        \hyperref[table:t1rq]{\textbf{RQ2}}};

        \node[box] (opensource) at (4.3,0)
        {Open-Source Repository\\
        Collection\\
        \hyperref[table:t1rq]{\textbf{RQ3}}};

        % ------------------------------------------------------------
        % Results of the three research streams
        % ------------------------------------------------------------

        \node[resultbox] (developerResult) at (-4.3,-2.2)
        {Logging Practices\\
        and Challenges};

        \node[resultbox] (literatureResult) at (0,-2.2)
        {Literature-Based\\
        Logging Guidelines};

        \node[resultbox] (opensourceResult) at (4.3,-2.2)
        {Repository-Based\\
        Logging Guidelines};

        % ------------------------------------------------------------
        % Open-source catalogue construction
        % ------------------------------------------------------------

        \node[box, text width=4.2cm] (catalogue) at (4.3,-4.8)
        {Open-Source Logging\\
        Guideline Catalogue\\
        Construction};

        % ------------------------------------------------------------
        % Thematic analysis
        % ------------------------------------------------------------

        \node[box, text width=4.2cm] (analysis) at (0,-7.2)
        {Thematic Analysis\\
        \hyperref[table:t1rq]{\textbf{RQ4}}};

        \node[resultbox, text width=4.2cm] (themes) at (0,-9.5)
        {Themes within the\\
        Logging Guideline Catalogue};

        % ------------------------------------------------------------
        % Discussion of findings
        % ------------------------------------------------------------

        \node[resultbox, text width=4.2cm] (discussion) at (0,-11.8)
        {Discussion of Findings};

        % ------------------------------------------------------------
        % Source -> result
        % ------------------------------------------------------------

        \draw[mainarrow]
        (developer) -- (developerResult);

        \draw[mainarrow]
        (literature) -- (literatureResult);

        \draw[mainarrow]
        (opensource) -- (opensourceResult);

        % ------------------------------------------------------------
        % Repository guidelines form the open-source catalogue
        % ------------------------------------------------------------

        \draw[mainarrow]
        (opensourceResult.south) -- (catalogue.north);

        % ------------------------------------------------------------
        % Literature provides the deductive codebook for RQ4
        % ------------------------------------------------------------

        \draw[mainarrow]
        (literatureResult.south)
        -- node[
            arrowlabel,
            pos=0.56,
            left
        ] {Literature-derived\\Deductive\\codebook}
        (analysis.north);

        % ------------------------------------------------------------
        % Open-source catalogue is analysed in RQ4
        % ------------------------------------------------------------

        \draw[mainarrow]
        (catalogue.south)
        -- ++(0,-1.75)
        -- (analysis.east);

        % ------------------------------------------------------------
        % Thematic analysis -> identified themes
        % ------------------------------------------------------------

        \draw[mainarrow]
        (analysis.south)
        -- node[
            arrowlabel,
            pos=0.50,
            left
        ] {Theme identification\\and refinement}
        (themes.north);

        % ------------------------------------------------------------
        % Themes contribute to the later discussion
        % ------------------------------------------------------------

        \draw[auxarrow]
        (themes.south) -- (discussion.north);

        % ------------------------------------------------------------
        % Developer study contributes independently to the discussion
        % ------------------------------------------------------------

        \draw[auxarrow]
        (developerResult.south)
        -- node[
            arrowlabel,
            midway,
            left
        ] {Developer perspective}
        (developerResult.south |- discussion.west)
        -- (discussion.west);

        \end{tikzpicture}%
    }

    \caption{Overview of the research design.}
    \label{fig:research_design}
\end{figure}

\section{Developer Study} \label{sec:meth_developerstudy}
The developer study investigates logging practices and challenges encountered during software development and maintenance.
It aims to provide practical insights from industry experience on logging to answer \hyperref[table:t1rq]{\textbf{RQ1}}.
The study focuses on how developers implement logging, analyse logs and use them for fault localisation.
For this purpose, a survey was conducted with software developers at an industry company. The company operates in the aviation field.
The following subsections describe the study design, participant recruitment, survey structure,
data collection procedure, and analysis of the collected responses in more detail.

\subsection{Study Design}
A survey was designed to collect information about developers' experiences with logging in real-world software projects.
The survey includes structured and free-text questions to collect quantitative and qualitative information about logging practices and challenges.
Structured questions are used to identify common patterns and frequently reported practices,
while free-text questions allow participants to explain their experiences in more detail.
This combination provides both comparable responses across participants and gives additional context.
The collected responses are used to identify common logging practices and recurring challenges encountered in practical software development.

\subsection{Participants and Recruitment}
The target group of the developer study consisted of software developers with
practical experience in software development and maintenance. Participants were
recruited from an industry company operating in the aviation field to obtain
insights from an industry development environment.
\\
Participation in the survey was voluntary. The questionnaire focused primarily on
participants’ experience with software logging and the difficulties encountered when
implementing, analysing and using logs for fault localisation. Detailed information
about individual projects or personally identifying information was not required for
the analysis.
\\
The survey dataset contained 26 recorded responses. Since not all participants
completed the entire questionnaire, the dataset contained both complete and partial
responses. Ten participants reached the final section of the survey, while other
participants responded only to earlier parts of the questionnaire.
\\
Partial responses were retained where a valid answer had been provided. Questions
that were not answered were excluded from the corresponding analysis and were not
interpreted as negative responses. This approach preserved useful information from
partially completed questionnaires while accounting for the number of valid
responses across individual survey questions.

\subsection{Survey and Data Collection}
The questionnaire was structured into three main parts covering logging implementation, log analysis and fault localisation.
Each part addressed different activities and challenges that developers may encounter when working with logs.
\\
The first part focused on logging implementation. Participants were asked about the amount of development effort spent on logging,
the data to log, the reasons for adding log statements and the tools or workarounds that have been used.
Further questions addressed how the relevance of logged information is assessed, how log levels are selected,
which log formats are used, and what difficulties arise when deciding what, where, and how to log.
\\
The second part focused on log analysis and examined how developers use logs during software development and maintenance.
Participants were asked for which purposes logs are analysed, how relevant log entries are identified
and which methods or tools are used to access and inspect log data.
Additional questions addressed working with unfamiliar logs, like structured, semi-structured and unstructured logs, collaboration with other developers
and difficulties that may occur during the analysis process.
\\
The third part investigated fault localisation and the role of logs in identifying the origin of software failures.
Participants were asked which additional information is required alongside log data and how important knowledge
about the software architecture, project context, development history and technical environment is for locating faults.
The questionnaire also considered the effort required for fault localisation, situations in which available log information
is insufficient and challenges caused by failures involving several system components.
\\
Different question formats were used throughout the questionnaire, including predefined response options,
multiple-selection questions, agreement scales, rankings, percentage estimates, and free-text responses.
The structured questions support comparison across participants, whereas free-text responses provide additional
context and allow participants to describe experiences that may not be covered by predefined answer options, such as
best practices and effective strategies for logging implementation, log analysis and fault localisation for each survey section.

\subsection{Data Analysis}

The survey responses were analysed using quantitative and qualitative methods.
Structured questions were evaluated using frequencies, percentages, agreement levels and rankings.
For multiple-selection questions, the number of selections for each response option was counted.
Agreement-scaled questions were used to identify general tendencies among the participants,
while percentage-based responses were used to estimate the effort required for different logging activities.
Because some questionnaires were only partially completed, the number of valid responses was considered separately for each question.
Missing answers were excluded from the corresponding analysis and were not treated as negative responses.
The free-text responses were lightly paraphrased for readability and conciseness while preserving their original meaning.
The paraphrased responses were manually checked against the original answers. No additional thematic coding or grouping was applied.
The free-text responses complement the structured responses by providing explanations and observations
that cannot be represented sufficiently through predefined answer options.
The findings from both forms of analysis are used to answer \hyperref[table:t1rq]{\textbf{RQ1}}
by identifying commonly reported logging practices and recurring challenges encountered during software development and maintenance.

\section{Collection of Logging Guidelines from Academic Literature}
\label{sec:meth_collection_literature}

To answer \hyperref[table:t1rq]{\textbf{RQ2}}, academic literature related to
software logging was collected and analysed. The primary aim was to identify
actionable recommendations for implementation and configuring
logging. Relevant publications that provided empirical, conceptual or tool-oriented
evidence were also retained where they contributed to understanding logging
practices, instrumentation or the analytical framework used in this thesis.
\\
Only identifiable actionable recommendations were treated as recommendation targets
for the later coverage comparison. Supporting observations, tool descriptions,
evaluation results and other non-prescriptive evidence were retained as literature
context but were not treated as independent logging guidelines.
\\
The actionable academic recommendations provide the reference set used for the
literature-based classification and the later comparison with the consolidated
open-source guideline catalogue.
\\
The following subsections describe the search strategy used to identify relevant
publications, the criteria applied to select the literature and the procedure used to
identify and extract actionable logging recommendations.

\subsection{Search Strategy}

Google Scholar was the primary platform for identifying relevant academic publications.
Several search terms were selected to cover different terminology used in research on software logging.
The following search terms were used:

\clearpage

\begin{itemize}
\item \textit{Logging Guideline}
\item \textit{Logging Practice}
\item \textit{Logging Survey}
\item \textit{What to Log}
\item \textit{Where to Log}
\item \textit{How to Log}
\end{itemize}

Using multiple search terms reduced the dependency on a single terminology, since logging recommendations
may be described as guidelines, practices, decisions or recommendations in different publications.
The search terms were intended to provide broad coverage of logging research rather than to correspond
one to one with the analytical dimensions used later in the study. Therefore, more specific topics
such as log levels and verbosity could also be identified through broader terms such as Logging Guideline or Logging Practice.
\\
The initial search focused on publications from 2020 to 2026.
Relevant publications identified through the search were additionally examined for references to earlier work
on software logging. When a cited publication provided relevant original research on logging practices or guidelines,
it was added to the literature set. This backward snowballing process was used to identify publications
that were not directly retrieved by the initial Google Scholar searches.
Regardless of the publication year, all relevant references identified through backward snowballing were considered for inclusion,
because older publications may contain important research that is still relevant to current logging practices.

\subsection{Selection Criteria}

Publications were selected based on their relevance to software logging
implementation, logging practices or the subsequent use of generated log data.
The literature set therefore contained both publications providing actionable
logging recommendations and publications providing relevant empirical, conceptual
or tool-oriented evidence.
\\
Publications containing actionable guidance were particularly relevant when they
addressed at least one of the four logging aspects considered in this thesis:
what information should be logged, where logging should occur, how logging should
be implemented or represented, or how log levels and verbosity should be selected.
\\
Empirical studies were also retained when they provided relevant evidence about
logging practices, instrumentation decisions or recurring logging problems.
Similarly, publications describing or evaluating logging tools, automated
instrumentation approaches, log parsers or related techniques could be retained as
supporting literature when they contributed to the conceptual background or
interpretation of software logging.
\\
A distinction was subsequently made between inclusion in the literature set and
inclusion in the academic recommendation set. A publication was treated as a source
of an actionable recommendation only when an identifiable statement prescribed,
recommended, or clearly demonstrated an action concerning the creation,
configuration or implementation of logging. Descriptive observations, tool
descriptions, evaluation results, model-performance measures and empirical findings
without a corresponding recommendation were not treated as independent coverage
targets.
\\
Publications were excluded when their relationship to logging was incidental and
they did not provide relevant evidence concerning logging practices,
instrumentation, log quality or the analytical context of the study. Duplicate
publications were also excluded.

\subsection{Guideline Extraction}

The selected publications were examined for statements containing recommendations
or guidelines concerning software logging. A statement was considered a logging
guideline when it explicitly described an action, recommendation or practice
concerning the creation or configuration of log information.
\\
Each identified guideline was extracted while preserving a reference to its original
publication. When a publication contained several distinct recommendations, these were
manually assigned to the corresponding logging category and recorded together in the catalogue for each publication.
The original context and source information were retained so that every guideline could be traced back to the
publication from which it originated.
\\
The extracted recommendations were manually summarised rather than copied word for
word. The purpose of this summarisation was to create concise and comparable guideline
statements while preserving the meaning of the original recommendation. No additional
recommendation, restriction or opposite recommendation was introduced during this
process.
\\
At this stage, semantically similar recommendations identified in different
publications were retained as separate source-level observations. This preserved
information about the independent occurrence of similar logging recommendations
across academic sources. The literature-derived recommendations were retained as a
separate reference set for construction of the initial deductive codebook and for the
later academic coverage comparison described in
Section~\ref{sec:meth_classification}.
\\
Each extracted guideline retained its original publication reference and was prepared
for subsequent use according to the main logging aspects
\textit{what to log}, \textit{where to log}, \textit{how to log}, and
\textit{log levels and verbosity}.

\section{Collection of Logging Guidance from Open-Source Repositories}  \label{sec:meth_collection_opensource}

To answer \hyperref[table:t1rq]{\textbf{RQ3}}, logging guidelines were collected from open-source software repositories hosted on GitHub.
Unlike academic literature, open-source projects do not necessarily document logging guidance in dedicated guideline documents.
Relevant recommendations may instead appear across different repository artefacts, such as README files, contribution guidelines,
project documentation, source files or configuration files.
Therefore, a multi-stage repository collection process was applied to identify, prepare and analyse relevant repository content.
The collected material was subsequently used to extract logging guidelines while preserving their connection to the source.
\\
The following subsections describe the repository search strategy, repository selection, collected repository characteristics
and the procedure used to identify and extract logging guidelines.

\subsection{Repository Search Strategy}

The repository search aimed to achieve broad coverage of logging guidance in open-source
software projects. Since logging recommendations may be expressed using different
terminology, several complementary GitHub search queries were used instead of relying
on a single search term.
\\
The initial query set was developed with the assistance of a large language model (GPT-5.4).
GPT-5.4 was used to suggest more precise terminology and combinations of search terms related to
different aspects of software logging for more precise results. The proposed queries were subsequently reviewed,
adapted and selected before being applied to GitHub. Therefore, GPT-5.4 supported the
formulation of candidate search queries, while the final search strategy and selection
of queries remained part of the research process.
\\
The final query set covered explicit logging guidelines as well as related topics
such as log content, structured logging, log levels, sensitive information,
observability and logging anti-patterns. Table~\ref{tab:github_queries}
summarises the queries used during the repository search.

\begin{table}[H]
    \centering

    \resizebox{1\textwidth}{!}{%
        \begin{tabularx}{\textwidth}{p{0.8cm} p{3.3cm} X}
            \hline

            \textbf{ID} & \textbf{Focus} & \textbf{Search Query} \\
            \hline

            Q1 &
            General logging guidelines &
            \texttt{(``logging guidelines'' OR ``logging best practices'' OR ``logging standards''
            OR ``logging conventions'' OR ``logging policy'') (path:/docs/ OR path:/.github/
            OR filename:README.md OR path:**/CONTRIBUTING.md)} \\
            \hline

            Q2 &
            What and when to log &
            \texttt{(``what to log'' OR ``log levels'' OR ``when to log'' OR ``structured logging''
            OR ``sensitive information'') (``logging'' OR ``application logs'')
            (path:/docs/ OR path:**/README.md)} \\
            \hline

            Q3 &
            Engineering standards and conventions &
            \texttt{(logging observability tracing telemetry)
            (``engineering standards'' OR handbook OR conventions OR practices OR guidelines)} \\
            \hline

            Q4 &
            Structured logging &
            \texttt{(``structured logging'' OR ``correlation id'' OR ``trace id'')
            (``best practices'' OR guidelines OR standards)} \\
            \hline

            Q5 &
            Security and sensitive information &
            \texttt{(logging OR logs)
            (``sensitive data'' OR PII OR GDPR OR credentials)
            (``guidelines'' OR ``best practices'')} \\
            \hline

            Q6 &
            Log levels and severity &
            \texttt{(``log level'' OR ERROR OR WARN OR INFO OR DEBUG)
            (``logging guidelines'' OR ``logging conventions'')} \\
            \hline

            Q7 &
            Repository popularity and programming language &
            \texttt{(``logging guidelines'' OR ``logging best practices'') OR
            (``logging guidelines'') language:Java OR
            (``structured logging'') language:Go} \\
            \hline

            Q8 &
            Markdown-based documentation &
            \texttt{(logging OR observability) path:*.md
            (guidelines OR practices OR standards)} \\
            \hline

            Q9 &
            Observability, telemetry and diagnostics &
            \texttt{(observability OR telemetry OR diagnostics OR tracing)
            (logging guidelines OR standards OR practices)} \\
            \hline

            Q10 &
            Logging anti-patterns &
            \texttt{(``logging smells'' OR ``bad logging practices'' OR ``logging anti-patterns'')} \\
            \hline

        \end{tabularx}%
    }

    \caption{GitHub search queries used for repository collection.}
    \label{tab:github_queries}
\end{table}

The queries were intended to complement each other by covering different terminology
and forms of logging guidance.
Specific search operators were used to focus on repositories and file types
in which guidelines are commonly documented.
\\
The results returned by the individual queries were combined into an initial search
dataset. Since the same repository or file could be returned by several queries,
duplicate handling and further preparation of the retrieved repositories were performed
in the subsequent repository selection step.

\subsection{Repository Retrieval}

After defining the search queries, an initial attempt was made to retrieve the search
results through the GitHub API. However, some of the predefined queries could not be
executed in their intended form because of restrictions affecting long and complex
search expressions containing several logical operators.
\\
Simplifying the queries would have changed the previously defined search strategy.
Therefore, the original query expressions were retained and executed through the
GitHub search interface instead.
\\
The searches were executed on 8 June 2026. For each
query, the first 10 search-result pages were collected, corresponding to a maximum of
approximately 100 returned results per query.
This boundary also reflected the GitHub search interface's limitations during data collection.
\\
The returned search-result pages were collected and processed using a Python script.
The script extracted the available repository names, file references and source URLs
from the search-result pages and stored these records for further processing. This
procedure was repeated for each of the defined search queries.
\\
The retrieved results from all queries were subsequently combined into an initial
dataset. At this stage, the complete set of retrieved records within the defined
search boundary was retained without removing duplicate repositories or files.
Cleaning and exclusion were performed only after the initial repository selection and
collection of repository characteristics.
\\
The use of a fixed number of search-result pages means that the collected dataset
represents the results available within the defined search boundary rather than all
repositories available on GitHub. This limitation is considered further in the
discussion of threats to validity.

\subsection{Repository Selection}

The results retrieved through the different search queries formed the initial
repository selection. The purpose of this stage was to preserve the complete set of
potentially relevant repositories and files identified through the predefined search
strategy.
\\
A repository was retained in the initial selection when the corresponding search
result contained potentially relevant logging-related material. No minimum threshold
based on the number of stars, commits, contributors or repository activity was applied
at this stage. Even repository creation date was considered later during the selection process.
This allowed both large and smaller open-source projects to be considered,
since relevant logging guidance may also occur in projects with lower levels of
development activity.
\\
Results returned by several different search queries were not removed at this point.
Similarly, repositories with overlapping characteristics or common authors were
initially retained. This ensured that metadata could first be collected for the complete
search result set before exclusion and duplicate handling were performed.
\\
The resulting initial selection therefore represents the complete repository dataset
obtained from the defined GitHub searches before data cleaning.

\subsection{Repository Characteristics}

Repository metadata were collected for the initial selection before duplicate
handling and other exclusion steps. The purpose of collecting these
characteristics was to provide contextual information about the projects in which
potentially relevant logging guidance occurred and to support later descriptive
analysis of the repository dataset.

The collected information included, where available:

\begin{itemize}
    \item repository and source reference,
    \item programming language,
    \item project type,
    \item application field or domain,
    \item number of commits,
    \item number of contributors or developers,
    \item repository activity,
    \item repository owner or author information.
\end{itemize}

Repository information directly available through GitHub, such as repository
reference, programming language, owner information, activity, commits, and
contributors were obtained from the corresponding repository information where
available.
\\
Project type and application field were determined from information available in the
repository documentation and description. Where AI-assisted (GPT-5.6 Sol) classification was
used to support the assignment of these characteristics, the assigned values were
reviewed against the corresponding repository information before being retained.
\\
Programming language, project type and application field were collected to describe
the technical and application context in which the logging guidance occurred. These
characteristics also support later descriptive analysis of whether particular
guidelines or themes occur across different types of projects or are concentrated
within particular parts of the collected dataset.
\\
The number of commits and contributors was interpreted as contextual information
about development activity and community involvement rather than as a measure of
repository or guideline quality. Repositories were therefore not excluded
because of few commits, contributors or other popularity indicators.
\\
Collecting the metadata before data cleaning additionally supported the comparison of
retrieved records and the identification of duplicate or overlapping repository
entries during the subsequent cleaning procedure.

\subsection{Data Cleaning and Exclusion}

After the initial repository selection and metadata collection, the dataset was
cleaned before guideline extraction.
\\
Since the same repository or file could be returned by several search queries,
duplicate search-result entries were identified and removed. Multiple occurrences
representing the same underlying repository or source material were consolidated
into a single entry.
\\
Repository owner and author information was additionally considered when identifying
overlapping results. However, repositories were not excluded because of
the same author or owner. Separate repositories were retained when projects 
or relevant material are represented. Entries were consolidated
only when duplicate or identical source material are referenced.
\\
Forked repositories were also examined during the cleaning process. Where a fork
contained the same relevant material as the corresponding original repository, the
original repository was retained, and the duplicated fork was excluded. When GitHub
search results displayed both an original repository and its forks, preference was
therefore given to the original repository to retain the source material
closest to its original project context.
\\
Additional duplicate and similarity detection was performed on the retrieved files
using cosine similarity. The similarity scores were used to identify files with
identical or highly similar textual content and to support the comparison of
potential duplicate material. A cosine similarity value of 1.0 indicated identical
representations of the compared file content. Files with almost identical content,
with a similarity value of around 0.9, were manually examined and duplicate
copies were removed before guideline extraction.
\\
Files with lower similarity values were not automatically removed based on their similarity score. Instead, similarity information was used to distinguish
between identical source material and files that contained related but potentially
different content. This prevented independently documented logging guidance from
being removed simply because its wording or structure was similar to another file.
\\
Repositories or referenced files that were no longer accessible on GitHub at the time
of data preparation were excluded. Information about excluded and unavailable
repositories was retained to keep the preparation procedure traceable.
\\
These cleaning steps concerned duplicate or overlapping source material within the
repository collection. Similar logging recommendations occurring independently in
different repositories were not removed at this stage. Such occurrences were retained
as independent source-level evidence and were considered separately during the later
construction and semantic consolidation of the logging guideline catalogue.
\\
After completion of the cleaning procedure, all remaining repositories and files
containing potentially relevant logging-related material were prepared for guideline
identification and extraction.

\subsection{Relevant Logging Content Identification and Extraction}
\label{sec:meth_content_extraction}

After data cleaning, the remaining relevant repository files were downloaded and
stored locally. The files were retained in their original file formats whenever
possible to preserve the source material and maintain traceability
to the repository and file from which the content originated.
\\
The collected files were analysed for logging recommendations concerning four
predefined logging dimensions:

\begin{itemize}
    \item \textit{what to log},
    \item \textit{where to log},
    \item \textit{how to log},
    \item \textit{log levels and verbosity}.
\end{itemize}

These four dimensions were used as high-level categories for identifying relevant
logging content within the repository files. At this stage, the purpose of the
classification was to locate potentially relevant source passages rather than to
assign the more detailed guideline codes introduced later during catalogue
construction.
\\
A passage was considered relevant when it contained an identifiable recommendation,
instruction or practice concerning the implementation or configuration of software
logging. This included both positive recommendations describing practices that should
be followed and negative recommendations describing practices that should be avoided.
Purely descriptive statements that mentioned logging without providing identifiable
guidance were not classified as logging recommendations.
\\
An AI-assisted (GPT-5.6 Sol) procedure was used to identify relevant logging content within the
collected repository files. Instead of asking the model to rewrite or summarise the
source material, the model was instructed to preserve the original content and add
XML tags around passages containing identifiable logging guidance.
\\
The following XML categories were used:

\begin{itemize}
    \item \texttt{<what\_to\_log>},
    \item \texttt{<where\_to\_log>},
    \item \texttt{<how\_to\_log>},
    \item \texttt{<log\_verbosity>}.
\end{itemize}

The original wording and structure of the source files were preserved as far as
possible. The model was instructed to add classification tags only around source
content containing identifiable logging guidance and not to generate, complement,
summarise or infer recommendations that were not explicitly supported by the
underlying material.
\\
Where a passage addressed more than one logging aspect, multiple corresponding
classifications could be retained. For example, a recommendation could specify both
which information should be recorded and how that information should be represented.
In such cases, the passage was not forced into only one of the four high-level
dimensions.
\\
Files for which no identifiable logging guidance was detected were retained without
additional XML classification tags. This prevented the procedure from requiring every
retrieved file to produce a guideline and reduced the risk of forcing unrelated or
purely descriptive content into one of the predefined categories.
\\
The tagging procedure was applied automatically to the collected files using
OpenAI GPT-5.6 Sol through the API. A predefined prompt was used consistently
across the processed dataset. The prompt defined what was considered relevant
logging guidance, specified the four classification categories and instructed
the model to preserve the original file content while adding XML tags only around
passages containing actionable logging recommendations or practices.
\\
The prompt used for this procedure was:

\begin{listing}[H]
\begin{minted}[linenos=false]{text}
I am providing you with a file that may contain logging guidelines.
Analyse the entire file and return the complete file with XML
classification tags added around passages that contain actual
logging guidelines. A logging guideline is a prescriptive or
recommended statement that tells developers what they should,
should not, must, may, or are expected to do when implementing
or configuring logging. This can also include clearly stated
best practices or recommendations expressed through examples.
Do not classify text merely because it discusses logging.
Only tag passages that contain actionable guidance, recommendations,
requirements, or explicitly demonstrated logging practices.
Classify each identified guideline using one or more
of the following categories:

* What to log <what_to_log>…</what_to_log>
* Where to log <where_to_log>…</where_to_log>
* How to log <how_to_log>…</how_to_log>
* Log verbosity <log_verbosity>…</log_verbosity>
Important:
* Keep the original wording exactly and dont change it at all.
* Only add tags if possible, else return file without any tags
\end{minted}
\caption{Prompt for logging content identification and XML classification}
\label{lst:prompt_xml_classification}
\end{listing}

After the AI-assisted (GPT-5.6 Sol) tagging procedure,
the generated files were checked for the
expected XML structure before their content was extracted. This technical integrity
check examined whether the expected opening and closing tags were present and whether
the generated output could be processed correctly by the subsequent extraction
script. Outputs containing malformed, incomplete or inconsistent XML structures were
identified and reprocessed or manually reviewed before being included in the
extraction stage.
\\
The technical XML check was distinguished from the validation of the content
classification itself. The accuracy of the AI-assisted (GPT-5.6 Sol) identification of relevant
logging passages was evaluated separately through manual validation as described in
Section~\ref{sec:meth_validation}. The validation included both files containing
identified logging guidance and files for which no guideline had been identified.
Including files without assigned tags made it possible to examine whether relevant
recommendations had potentially been missed during the automated identification
process.
\\
After the tagging procedure had been completed, the XML-marked passages were
extracted automatically using a Python script and transferred into a structured
open-source logging-content corpus. The extraction procedure identified the content
contained within the four predefined XML categories and stored the corresponding
passages in a structured representation.
\\
Multiple distinct passages originating from the same source file were retained as
separate entries where applicable. Similarly, comparable recommendations occurring
independently in different repositories were retained separately rather than being
combined during this stage.
\\
Each extracted passage remained linked to its original repository and source file.
The assigned high-level logging category was also retained to support the subsequent
guideline formulation and classification.
\\
The resulting dataset therefore preserved the relationship between the extracted
logging content, its classification and its original repository context. This
traceability enabled the extracted passages to be compared with their source
material during subsequent normalisation and validation.
\\
At this stage, the extracted passages represented source-level logging content and
were not yet treated as normalised logging guidelines. No semantic consolidation of
similar recommendations was performed. Consequently, similar recommendations
documented independently in different repositories remained separate source-level
observations.
\\
The subsequent transformation of the extracted source content into concise normalised
guideline statements, the manual review of these statements, their comparison with
the academic reference set, detailed classification and semantic consolidation
are described in Section~\ref{sec:meth_classification}.

\section{Logging Guideline Catalogue Construction and Classification}
\label{sec:meth_classification}

The preceding sections describe the collection of logging guidelines from academic
literature and the identification of relevant logging content within open-source
repositories. While the academic literature already provided explicitly formulated
logging recommendations, the open-source material differed considerably in wording,
structure and level of detail. A separate catalogue construction process was therefore
applied to transform the collected material into a consistent and analysable set of
logging guidelines.
\\
The objective of this stage was to create a traceable open-source guideline catalogue while
preserving the relationship between each guideline and its source. The
process consisted of contextual content tagging, guideline formulation and normalisation,
manual review, classification using an initial literature-based codebook, inductive
extension of the codebook, semantic consolidation and final validation. The retained
academic recommendations were used to derive the deductive starting structure and were
compared with the consolidated catalogue after catalogue construction.
\\
The resulting catalogue forms the basis for the thematic analysis described in
Section~\ref{sec:meth_thematic_analysis}.

\subsection{Contextual Content Tagging}

Following extraction of the XML-classified source passages, an additional contextual
tagging step was applied before the passages were transformed into normalised logging
guidelines. The purpose of this step was to retain a short description of the main
topic or topics represented by each extracted passage.
\\
For each extracted entry, the source passage was retained in the
\textit{Content} field. An AI-assisted procedure was then used to create a corresponding
\textit{Content Tag}. The Content Tag was used as contextual metadata to support later
guideline formulation and validation rather than as the final thematic classification
of the guideline.
\\
The prompt used for this procedure was:

\begin{listing}[H]
\begin{minted}[linenos=false]{text}
I am providing an Excel file containing extracted logging content.
Please analyse the Content column and add a new column called Content Tag.
For each row, create a short summary of the main topic
or topics discussed in the content.
\end{minted}
\caption{Prompt for contextual content tagging}
\label{lst:prompt_content_tagging}
\end{listing}

The original Content remained available alongside the Content Tag so that the technical
context could be checked when a short tag was ambiguous. During manual investigation,
contextual expressions that could lead to misleading interpretations were checked
against the source passage and corrected where necessary.

\subsection{Guideline Formulation and Normalization}

The logging content extracted from the open-source repository files was initially
retained in its original wording. Since the extracted passages differed in length,
writing style and level of detail, they were transformed into guideline statements
before further analysis.
\\
An AI-assisted (GPT-5.6 Sol) procedure supported this normalisation step.
The original Content and its related Content Tag were analysed together to identify
the logging recommendation contained in the source, and turn it into a clear and
independently understandable logging guideline. The aim was to express recommendations
that consistently appeared in different forms so that they could be compared
within the catalogue. No additional logging advice was added, and the final guideline
was based only on the information available in the original context.
\\
Guidelines could also be derived from clearly demonstrated logging practices in code,
configuration, examples or tables where the source material sufficiently supported the recommendation. Information that was not supported by the source was not added.
The normalised guideline was formulated so that it could be understood independently
without referring back to the Content field.
\\
Where no meaningful logging guideline could be derived from the extracted material,
the entry was assigned \texttt{NO\_GUIDELINE}. Such entries were retained separately
and were not assigned themes during the subsequent classification.
These entries were also manually reviewed to confirm that no relevant
logging guideline had been overlooked.
\\
The original Content and Content Tag were retained alongside the normalised guideline.
Consequently, every transformed guideline could subsequently be compared with the
source material from which it had been derived.
\\
The prompt used for the procedure was:

\begin{listing}[H]
\begin{minted}[linenos=false]{text}
Analyse the Content column and content tag, and add a new Guideline column.
Transform each row into a clear, concise, self-contained and stand
alone as logging guideline, where content is not reference in the guideline.
Guidelines may also  be derived from clearly demonstrated practices
in code, configuration, examples, or tables.
Preserve the source meaning and do not add  unsupported information.
Use NO_GUIDELINE only  when no meaningful guideline can be derived or
content tag include no guideline.
Add all NO_GUIDELINE rows to a separate worksheet.
\end{minted}
\caption{Prompt for guideline formulation and normalization}
\label{lst:prompt_guideline_normalization}
\end{listing}

\subsection{Manual Review of Normalised Guidelines}

Following normalisation, the generated guideline statements remained linked to their
corresponding source passages and Content Tags so that the transformation
could be manually validated.
\\
A random sample of approximately 10\% of the normalised entries was manually
compared with the corresponding Content and Content Tag. The review considered whether
the normalised guideline represented the recommendation contained in the source,
whether relevant restrictions or conditions had been preserved, whether the guideline
was sufficiently self-contained and whether additional unsupported meaning had been
introduced during normalisation.
\\
Additional manual investigation was performed for entries identified during the
processing procedure as unclear or inconsistent. Where recurring problems were
identified, related entries were checked and corrected rather than changing only the
individual example in which the problem had first been observed.
\\
Entries classified as \texttt{NO\_GUIDELINE} were also checked to ensure that no
meaningful recommendation was excluded only because of the automated transformation.
\\
After this review, the resulting entries formed the source-level open-source guideline
catalogue used to address \hyperref[table:t1rq]{\textbf{RQ3}}. Each guideline remained linked to its original
passage, repository and source file during the repository collection process.

\subsection{Relationship Between Academic and Open-Source Guidelines}

The academic and open-source guideline sets were retained as analytically distinguishable
source sets. The open-source recommendations formed the source-level catalogue that was
normalised, thematically classified and semantically consolidated. The academic
recommendations were retained with their publication references as the literature-based
reference set, while descriptive literature remained available as supporting evidence.
\\
The academic literature contributed to the catalogue analysis in two ways. First,
classifications identified in previous logging research were used to construct the
initial deductive codebook described in the following subsection. Second, after semantic
consolidation of the open-source recommendations, the retained academic recommendations
were manually compared with the resulting canonical guidelines to determine whether each
academic recommendation was represented in the final consolidated open-source catalogue.
\\
The comparison was performed at the consolidated guideline level rather than separately
for every repeated source occurrence. This prevented a recommendation documented in
several open-source sources from being counted repeatedly when evaluating its relationship
to the academic recommendations.
\\
The academic and open-source provenance remained retained throughout this comparison.
The methodological procedure for the comparison is described later in the thematic
analysis, while the actual academic coverage identified for the final consolidated
guidelines is reported in the results chapter.

\subsection{Initial Literature-Based Codebook}

An initial deductive codebook was constructed from classifications and concepts
identified in previous software logging research. The purpose of this codebook was
to provide an academically grounded starting structure for the classification of the
collected guidelines rather than to define an exhaustive set of possible themes.
\\
The four high-level categories used during the earlier extraction process were
retained:

\begin{itemize}
\item \textit{what to log},
\item \textit{where to log},
\item \textit{how to log},
\item \textit{log levels and verbosity}.
\end{itemize}

For \textit{what to log}, the initial codes were based on the classification reported
in \textit{Industry Practices and Event Logging: Assessment of a Critical Software
Development Process}~\cite{peccia2015industrypracticesandeventlogging}. The following
codes were used:

\begin{itemize}
\item \textit{state dump},
\item \textit{execution tracing},
\item \textit{event reporting}.
\end{itemize}

For \textit{where to log}, the initial codes were derived from
\textit{Where Do Developers Log? An Empirical Study on Logging Practices in Industry}
~\cite{fu2014wheredodeveloperslog}. The following classifications were included:

\begin{itemize}
\item \textit{assertion-check logging},
\item \textit{return-value-check logging},
\item \textit{exception logging},
\item \textit{logic-branch logging},
\item \textit{observing-point logging}.
\end{itemize}

For \textit{how to log}, the initial classification was based on
\textit{A Survey of Software Log Instrumentation}
~\cite{chen2021asurveyofsoftwareloginstrumentation}. The survey distinguishes three
main stages of log instrumentation:
\clearpage
\begin{itemize}
\item \textit{logging approach},
\item \textit{logging utility integration},
\item \textit{logging code composition}.
\end{itemize}

For \textit{log levels and verbosity}, the initial codes were based on the verbosity
levels discussed by Gholamian and Ward \cite{gholamian2021whatdistributedsystemssay}.
The paper refers to the six default Log4j verbosity levels:
\begin{itemize}
\item \textit{FATAL},
\item \textit{ERROR},
\item \textit{WARN},
\item \textit{INFO},
\item \textit{DEBUG},
\item \textit{TRACE}.
\end{itemize}

The resulting initial codebook therefore consisted of literature-derived codes within
each of the four predefined logging categories. Guidelines that could not be
meaningfully represented by one of these codes were not forced into the existing
classification. Instead, they were marked for subsequent investigation.
\\
The codebook was based on classifications identified in the academic literature.
For each code, the corresponding literature was used to determine its meaning and
to guide its application during classification.

\subsection{Inductive Codebook Extension and Final Classification}

The literature-derived codes provided the initial classification structure for the
guideline catalogue. However, not all guidelines could be represented adequately by
these initial codes. Additional themes were therefore introduced based on concepts
identified within the collected guideline content.
\\
The normalised Guideline column was analysed using the literature-derived codebook
first. The four high-level logging categories remained \textit{what to log},
\textit{where to log}, \textit{how to log} and \textit{log verbosity}. Where a
guideline could be represented by one of the existing deductive themes, the
corresponding theme was assigned. If both a broad and a more specific deductive theme
were applicable, the most specific applicable theme was selected.
\\
Where none of the existing deductive themes adequately represented the guideline, a
concise additional theme was derived inductively from the guideline content. For each
assignment, the origin of the theme was retained in the field \textit{Theme Basis} as
either \textit{Deductive} or \textit{Inductive}. Guidelines assigned an inductive theme
were additionally collected in a separate review worksheet so that the newly derived
themes could be inspected and later harmonised.
\\
Entries classified as \texttt{NO\_GUIDELINE} were not assigned a Theme or Theme Basis.
\\
The prompt used for this classification was:

\begin{listing}[H]
\begin{minted}[linenos=false]{text}
Analyse the Guideline column and add Theme and Theme Basis and Leave
out for NO_GUIDELINES
Use this codebook first:
What to log: State dump, Execution tracing, Event reporting
Where to log: Assertion-check logging, Return-value-check logging,
Exception logging, Logic-branch logging, Observing-point logging
How to log: Logging approach, Logging utility integration,
Logging code composition
Log verbosity: FATAL, ERROR, WARN, INFO, DEBUG, TRACE
Check all deductive themes first, then if a guideline fits both a broad
and a more specific deductive theme, assign the most specific applicable theme,
else when no existing theme fits, create a concise
new inductive theme from the guideline.
Set Theme Basis to Deductive or Inductive.
Put all inductively themed rows in a separate worksheet called Inductive Review.
Leave the NO_GUIDELINES empty for Theme and Theme Basis
\end{minted}
\caption{Prompt for deductive-inductive theme classification}
\label{lst:prompt_theme_classification}
\end{listing}

The theme assignment was based on the semantic meaning of the normalised
guideline. Each guideline was assigned to one final high-level logging
dimension and one corresponding theme. Where a guideline or its underlying
source content had previously been associated with more than one dimension,
the original Content was consulted, and the guideline was assigned to the
dimension that most specifically represented its primary recommendation.
The original Content, Content Tag and source-level classifications remained
available for contextual interpretation and traceability. The Content Tag
itself was not automatically treated as a Theme.
\\
After the initial classification, inductively generated themes were reviewed across
the catalogue. Overlapping inductive themes representing the same concept were merged, and their names standardised, while distinct concepts were retained separately. The
existing literature-derived themes retained their deductive basis.
\\
The resulting classification therefore combined literature-derived deductive themes
with additional themes derived from the collected guideline content. The analytical
structure consisted of the high-level logging dimension, the assigned theme and the
individual guideline. Guidelines assigned to the same theme were not assumed to be
semantically equivalent, since the guideline itself retained the specific recommendation
and context.

\subsection{Guideline Consolidation and Traceability}

Following classification, the source-level catalogue was examined for semantically
equivalent guideline statements. This step was separate from the removal of duplicate
repository material performed during data cleaning. Recommendations independently
documented in different sources were treated as separate source occurrences even when
their meaning was equivalent.
\\
Guidelines were considered candidates for consolidation when they described the same
underlying normative recommendation. Equivalence was assessed by considering the
recommended action, the logging concern to which the action referred and any conditions
or restrictions attached to the recommendation. Guidelines addressing a similar topic
but containing a distinct recommendation, condition or restriction were retained
separately. Uncertain cases were also kept separate.
\\
When semantically equivalent source-level guidelines occurred under different
dimensions or themes, their original Content was compared before consolidation
to determine whether the apparent equivalence resulted from normalisation or
represented the same underlying recommendation. Where consolidation was
appropriate, the resulting canonical guideline was assigned to the single
dimension and theme that most specifically represented its primary
recommendation. The original source-level dimension and theme assignments were
retained in the consolidation mapping for traceability. Where the underlying
recommendations were meaningfully different, or the equivalence remained
uncertain, the guidelines were retained separately.
\\
Due to the number of guidelines contained in the catalogue, AI-assisted semantic
comparison using GPT-5.6 Sol was used to support the identification and consolidation
of semantically equivalent guidelines. The model was instructed not to remove unique
recommendations or conditions and to preserve the mapping from every original guideline
to the corresponding consolidated guideline.
\\
A separate code-review indicator was used to identify source Content containing
technical material that could later be displayed as an illustrative example for a
guideline. This included source code, configuration, use-case or implementation examples,
example input and program output. In the working catalogue, this information was retained
through the \textit{Code Involved} field. The term \textit{Code Review} in the
consolidation prompt referred to this example-identification step and did not indicate
a methodological validation status.
\\
For entries identified in this way, the associated source Content was preserved as a
separate \textit{Example Reason} linked to the consolidated guideline. These examples
were not merged, removed or generalised during consolidation, even when their underlying
guideline was consolidated.
\\
The prompt used for the consolidation procedure was:

\begin{listing}[H]
\begin{minted}[linenos=false]{text}
Merge semantically equivalent guidelines without losing any unique
recommendation or condition. Keep uncertain cases separate.
For every row with Code Review = Yes, preserve its Content as a
separate Example Reason linked to the consolidated guideline and do not
merge, remove, or generalise these examples during consolidation.
Create consolidated Guidelines and consolidation mapping
Then update the codebook: keep existing deductive themes as
Deductive, merge overlapping inductive themes, standardise names,
and keep distinct concepts separate.
create updated codebook and theme mapping
Ensure every original guideline is represented exactly once
and summarise identical guidelines together. Provide a sheet
with all summarized guidelines, their content, type, theme for manual
investigation later
\end{minted}
\caption{Prompt for guideline consolidation and codebook refinement}
\label{lst:prompt_guideline_consolidation}
\end{listing}

Guidelines confirmed as semantically equivalent were mapped to a common canonical
guideline. The original source-level guideline records were retained to
preserve traceability. Each original guideline was required to be represented exactly
once in the consolidation mapping.
\\
The consolidation stage also updated the codebook and theme mapping. Existing deductive
themes remained deductive. Overlapping inductive themes were merged and standardised
where they represented the same underlying concept, while distinct concepts were kept
separate.
\\
This approach reduced repetition in the final presentation of the catalogue while
preserving provenance, contextual examples and the independent source occurrences
associated with each consolidated recommendation. Exact duplicates originating from
repeated retrieval of the same underlying source remained handled by the earlier
data-cleaning stage and were not counted as independent source occurrences.

\subsection{Validation of the Guideline Catalogue}
\label{sec:meth_validation}

The AI-assisted (GPT-5.6 Sol) catalogue construction involved several processing stages
at which identification, transformation or classification errors could occur. Validation
was therefore performed at multiple stages of the catalogue construction process.
\\
First, the XML-tagged repository files were checked for technical integrity. The check
examined whether the expected opening and closing tags were present and whether the
output could be processed correctly by the extraction script. Malformed, incomplete or
inconsistent XML outputs were identified and reprocessed or manually investigated.
\\
The semantic content of the extracted material was additionally reviewed during the
subsequent preparation steps. Manual investigation was used to identify cases in which
a contextual expression had been interpreted incorrectly, or the high-level XML category
did not adequately represent the source content. Examples included log-verbosity cases,
unstructured log-message content, formatting-related recommendations and HTTP-related
content. In particular, HTTP-related entries were investigated because occurrences of
HTTP could represent hyperlinks or references rather than an HTTP-specific logging
recommendation. Where recurring issues were identified, corresponding entries were
reviewed and corrected throughout the affected part of the dataset.
\\
Second, the transformation from Content and Content Tag to normalised guidelines was
validated manually. A random sample of 234 entries, corresponding to approximately 10\% of the dataset at this stage, was compared with the associated source Content
and Content Tag. The review examined whether the guideline preserved the source meaning,
retained relevant conditions and restrictions, remained self-contained and did not
introduce unsupported information. Cases requiring additional context were corrected
using the original Content as the reference.
\\
Entries assigned \texttt{NO\_GUIDELINE} were additionally reviewed to ensure that the
classification did not exclude meaningful logging recommendations. Content that only
contained references, links or descriptive material without a supported logging
recommendation remained excluded from thematic classification.
\\
Third, the final theme classification was validated using the deductive–inductive
codebook. Validation was performed across the individual themes rather than relying on
one undifferentiated sample from the complete dataset. Larger themes were checked using
manual samples, while themes with comparatively few assigned guidelines were inspected
more closely because individual classification errors would have a larger proportional
influence on their reported distribution.
\\
The guidelines initially assigned to \textit{Logging approach} were re-examined against 
the complete final codebook and, where necessary, their original Content to determine 
whether a more specific existing theme represented the recommendation more accurately. 
This review aimed to determine whether the guidelines fell within the theme's conceptual scope. 
Guidelines within such a broad theme did not need to be semantically equivalent, 
since each guideline retained its specific scenario and recommendation.
In contrast, the theme served as a higher-level classification.
\\
During this validation, incorrectly assigned log levels, event-reporting cases and
logging-approach cases were corrected where identified. When a repeated classification
pattern was detected, related entries were investigated rather than correcting only the
initial sampled row. The final theme distribution was calculated after these corrections.
\\
Finally, the semantic consolidation and traceability of the catalogue were checked.
Consolidated guidelines were reviewed to determine whether semantically equivalent
recommendations had been combined without removing a unique condition, restriction or
recommendation. Where equivalence remained uncertain, the guidelines were retained
separately.
\\
Where manual validation showed that a consolidated group contained source occurrences
representing different recommendations, the affected occurrences were separated and,
where an equivalent canonical guideline already existed elsewhere in the catalogue, 
reassigned to that guideline rather than creating a duplicate recommendation.
A new canonical guideline was retained only when the recommendation was not adequately
represented by an existing guideline.
\\
An integrity check was also performed on the consolidation mapping. Every retained
source-level guideline was required to be represented exactly once by a consolidated
canonical guideline, and every canonical guideline remained linked to the source-level
guideline or guidelines from which it had been derived. Example Reasons originating
from code, configuration, usage, implementation or output examples remained separately
linked to the corresponding consolidated guideline.
\\
The final validation additionally checked that each canonical guideline was
assigned to one final high-level logging dimension and corresponding theme.
Where consolidation had combined equivalent source occurrences with different
earlier classifications, the most specific classification supported by the
underlying recommendation was retained as the final assignment, while the
original classifications remained traceable in the mapping.
\\
The resulting catalogue therefore preserves the relationship between the 
source content, contextual Content Tag, normalised source-level guideline, assigned
theme and consolidated canonical guideline. This structure supports later manual
investigation while preserving the provenance of the final catalogue.

\section{Thematic Analysis}
\label{sec:meth_thematic_analysis}

The final classified and consolidated logging guideline catalogue was used as the
basis for answering \hyperref[table:t1rq]{\textbf{RQ4}}:
\textit{What themes can be identified within the logging guideline catalogue?}
\\
A hybrid deductive–inductive thematic classification was applied. The deductive
component originated from the literature-derived initial codebook, while the inductive
component allowed additional themes to be introduced for guideline concepts that could
not be represented adequately by the initial classification.
\\
The thematic analysis considered both the final themes and the individual guidelines
assigned to them. The four predefined logging aspects \textit{what to log},
\textit{where to log}, \textit{how to log}, and \textit{log levels and verbosity}
were retained as the high-level analytical dimensions. Within these dimensions, the
final deductive and inductive themes represented recurring areas of logging guidance.
\\
The individual normalised, or consolidated guideline remained the more specific unit of
actionable knowledge. Consequently, guidelines assigned to one broad theme were not
interpreted as identical recommendations. The themes provided an organisational and
analytical structure through which recurring logging concerns could be examined across
the catalogue.

\subsection{Theme Identification}

Theme identification followed the hybrid deductive–inductive classification procedure
described in Section~\ref{sec:meth_classification}. Each guideline was first compared
with the literature-derived deductive themes. Where an existing deductive theme
adequately represented the recommendation, that theme was assigned. Where no existing
theme was sufficient, an inductive theme was created based on the semantic meaning of
the guideline.
\\
The themes were reviewed iteratively across the catalogue. Overlapping inductive themes
representing the same underlying concept were merged, and their names standardised,
while themes representing distinct logging concepts remained separate. Existing
deductive themes retained their literature-derived basis.
\\
Theme development was based on the meaning of the underlying guidelines rather than
only on similarity in terminology. The source passages, Content Tags and
normalised guidelines were consulted where necessary to preserve the context of a
recommendation during classification.
\\
The resulting structure consisted of the four high-level logging dimensions, the final
deductive or inductive Theme assigned within the corresponding dimension, and the
specific guideline. This structure allowed the themes to be analysed descriptively
while retaining the detailed recommendations contained in the catalogue.

\subsection{Theme Distribution}

After the themes had been identified, their distribution within the guideline
catalogue was analysed descriptively. This analysis considered both the number of
canonical guidelines associated with each theme and the number of source-level
occurrences contributing to those guidelines.
\\
The distinction between canonical guidelines and source occurrences was maintained
throughout this analysis. A canonical guideline represents a semantically consolidated
recommendation, whereas a source occurrence represents an independent appearance of
that recommendation within an open-source repository.
\\
For each identified theme, the following information was considered where applicable:

\begin{itemize}
\item number of canonical guidelines,
\item number of source-level guideline occurrences,
\item number of contributing open-source repositories,
\item distribution across the four high-level logging categories.
\end{itemize}

The number of occurrences was used as a descriptive indication of how frequently a
theme appeared within the collected corpus. Frequency was not interpreted as a
measure of guideline quality, correctness or importance. A frequently occurring
recommendation may indicate that different sources repeatedly address a concern, while a less frequent recommendation may still represent a relevant or
specialised logging practice.
\\
Where appropriate, repository metadata such as programming language, project type
or application field was used as additional descriptive context. Such comparisons
were used to explore whether particular themes occurred across different types of
projects or were concentrated within specific parts of the collected dataset.

\subsection{Coverage of Academic Recommendations}

After semantic consolidation, the retained actionable academic recommendation
targets were manually compared with the canonical open-source guidelines. The
purpose of this comparison was to determine the extent to which recommendations
identified in academic literature were represented in the final consolidated
open-source catalogue.
\\
The unit of comparison was the individual actionable recommendation rather than the
publication as a whole. A publication could therefore contribute several distinct
recommendation targets, while publications retained only for empirical, conceptual
or tool-oriented context did not contribute a coverage target.
\\
The comparison was performed against the consolidated catalogue rather than
separately against each repeated source occurrence. Where equivalent guidance
appeared independently in several repositories, the corresponding canonical
guideline was used for the coverage assessment, while the source-level occurrences
remained available for the separate frequency analysis.
\\
The comparison considered the semantic meaning of the academic recommendation and
the corresponding canonical guideline, including its recommended action, relevant
conditions, scope and restrictions.
\\
An academic recommendation was classified as fully covered when one or
more canonical guidelines expressed the same underlying recommendation while
preserving the relevant meaning and conditions.
\\
A recommendation was classified as partially covered when the principal
logging concern was represented in the catalogue but one or more relevant aspects,
conditions or sub-recommendations were not represented sufficiently.
\\
A recommendation was classified as not identified when no canonical
guideline in the collected open-source catalogue sufficiently represented the
underlying recommendation. Thematic similarity alone was not considered sufficient
for guideline-level coverage.
\\
The retained provenance of both datasets allowed each academic recommendation to
remain linked to its publication and each matching recommendation to remain linked
to the corresponding canonical guideline or guidelines. The resulting coverage
classification is reported in the Results chapter.
\\
The comparison is limited to the selected academic recommendation set and the
collected open-source corpus. A partial or missing match therefore does not
demonstrate that a recommendation is absent from software logging research or
open-source software in general.
